% ===================================================================
%  TABEL Nr. 11 — Linn ja selle monumendid. --- Tulekahju.
%  Traduction 2026 d'apres texto/io, controlee sur texto/fr et
%  texto/fi. Consigne : voir texto/et/10-tabel-01.tex.
%
%  LE TABLEAU 11 ROMANCHE AVAIT TROUVE UN COUPLE QUI SE PARTAGEAIT DE
%  LA MEME FACON DANS DEUX LANGUES SANS RAPPORT ; L'ESTONIEN EST LA
%  TROISIEME, ET IL CASSE LE COUPLE.
%
%  Le luxembourgeois : « Spuerkeess » la caisse d'epargne, compose a
%    la maison ; « Pompjeeën » les pompiers, emprunte tel quel.
%  Le romanche : « banca da spargn » compose ; « pumpiers » emprunte.
%  L'ESTONIEN : « hoiukassa » la caisse d'epargne, compose ; et
%    « tuletõrjujad » les pompiers, LES REPOUSSEURS-DE-FEU, compose
%    lui aussi, de « tuli » le feu et « tõrjuma » repousser. Pas
%    l'ombre d'un emprunt.
%
%  DEUX LANGUES DISAIENT LA MEME CHOSE, LA TROISIEME DIT LE
%  CONTRAIRE : IL FAUT DONC QUE LE CRITERE AIT ETE MAL MESURE. On
%  avait ecrit que le corps de pompiers « arrive du dehors deja
%  constitue » et apporte son nom. C'est vrai du Luxembourg et des
%  Grisons ; ce n'est pas vrai de l'Estonie.
%
%  Les corps de pompiers volontaires estoniens ont ete fondes dans
%  les annees 1860-1880 par les societes du reveil national — celles-
%  la memes qui, aux memes annees, fabriquaient des mots. Le MODELE
%  venait d'Allemagne ; LA FONDATION etait estonienne, et elle s'est
%  nommee elle-meme.
%
%  LE CRITERE TIENT DONC, MAIS IL FAUT LIRE « FONDER » AU SENS
%  ETROIT : non pas qui a invente l'institution, mais QUI A SIGNE LES
%  STATUTS DE CELLE-CI, ICI. Les pompiers du Luxembourg ont ete
%  installes ; ceux de l'Estonie ont ete fondes. Une idee importee et
%  une association locale ne se nomment pas de la meme facon, et la
%  serie confondait les deux depuis le tableau 11 luxembourgeois.
%
%  C'est la deuxieme fois que la colonne estonienne corrige une regle
%  au lieu de la confirmer : le tableau 2 avait deja montre qu'une
%  langue planifiee refuse le mot sans refuser la chose.
% ===================================================================

%%K t11-apar-1 apar
\VUsaut{2.0mm}
\VUcentre{13.2pt}{90}{KOLMAS SEERIA}
\VUsaut{3.0mm}
\VUfilet{16mm}
\VUsaut{5.6mm}
\VUtitre{15.0pt}{60}{TABEL Nr. 11}
\VUsaut{1.4mm}
\VUpk{11.4pt}{40}{Linn ja selle monumendid. --- Tulekahju.}
\VUsaut{2.2mm}

%%K t11-tit-1 sub
\VUpk{10.2pt}{40}{Eeslinn.}
\VUsaut{3.0mm}

%%K t11-c1-01-1 p
1. --- Ma elan suures linnas, mis asub 100 kilomeetri kaugusel ookeanist ja mis on
siiski esimese järgu \VUgras{sadam} \textsuperscript{(1)}. Jõgi, mis seda ääristab, on
400 meetrit lai ja moodustab väga ilusa reidi.

%%K t11-c1-02-1 p
2. --- Kogu linna osa, mis on \VUgras{kaidel} \textsuperscript{(2)}, näib täiesti
merelisena ja tööstuslikuna ning moodustab \VUgras{eeslinna} \textsuperscript{(3)}, mida
asustavad ainult meremehed ja töölised. Need töötavad \VUgras{vabrikutes}
\textsuperscript{(4)} ja \VUgras{gaasitehases} \textsuperscript{(5)}, mille kõrge
\VUgras{korsten} \textsuperscript{(6)} ja \VUgras{gaasimahuti} paistavad väga kaugele.

%%K t11-c1-03-1 p
3. --- Keskajal oli linn kindlustatud, ja leidub veel mõned jäänused selle müüridest,
eriti \VUgras{värav} \textsuperscript{(8)}, mis kuulub vanale \VUgras{kindlustatud
lossile} \textsuperscript{(7)} ja mida külgnevad selle kaks \VUgras{torni}
\textsuperscript{(9)}, mida praegu kasutatakse \VUgras{vanglana}.

%%K t11-c1-04-1 p
4. --- Kogu selles \VUgras{linnaosas}, mis näib väga vana, on ainult üks ehitis
suhteliselt uus: see on \VUgras{toll} \textsuperscript{(10)}. See asub kai ja
\VUgras{tänavakese} \textsuperscript{(11)} vahel, mis viib vana \VUgras{raekoja}
\textsuperscript{(12)} juurde. Seda viimast monumenti tuntakse kergesti ära hiiglasliku
\VUgras{kellatorni} \textsuperscript{(13)} järgi, mis valitseb vana linna.

%%K t11-c1-05-1 p
5. --- Tänava teisel pool \VUgras{seisab} hoone, mis on ehitatud samas stiilis nagu
toll: see on \VUgras{börs} \textsuperscript{(14)}. Üks selle fassaadidest vaatab
väiksele platsile, kus on \VUgras{prefektuur} \textsuperscript{(15)}. Tõepoolest, meie
linn, olemata pealinn, on siiski tähtis departemangu keskus.

%%K t11-tit-2 sub
\VUsaut{5.0mm}
\VUpk{10.2pt}{40}{Linna kõrgeim osa.}
\VUsaut{3.4mm}

%%K t11-c1-06-1 p
6. --- Garnison, üsna arvukas, on majutatud avaratesse väga tervislikesse
\VUgras{kasarmutesse} \textsuperscript{(16)}; need ulatuvad \VUgras{linnapargi}
\textsuperscript{(17)} võreni. \VUgras{Puiestee} \textsuperscript{(19)}, mis viib sellesse
parki, möödub \VUgras{hoiukassa} \textsuperscript{(18)} tagant ja lõpeb
\VUgras{katedraali} \textsuperscript{(20)} platsil.

%%K t11-c1-07-1 p
7. --- Kõik need monumendid laiuvad amfiteatrina väiksel künkal, kus on ka meie avar
\VUgras{gümnaasium} \textsuperscript{(21)}.

%%K t11-c1-07-2 p
Kui minna alla veidi kaldu teed mööda, mis ühineb Suure tänavaga, jõutakse
\VUgras{kohtumajani} \textsuperscript{(22)}, mille ees laiub meeldiv \VUgras{avalik aed}
\textsuperscript{(26)}. See \VUgras{aiaplats} ei ole väga suur, kuid see teeb keset linna
rohelise ja jaheda oaasi. Seepärast käivad seal palju linlased, kellele meeldib tulla
istuma \VUgras{kohviku} \textsuperscript{(25)} varikatuse alla, et kuulata
sõdurmuusikuid, kes neljapäeviti ja pühapäeviti mängivad \VUgras{kioskis}
\textsuperscript{(27)}, mis on pandud murude ja lillepeenarde vahele.

%%K t11-c1-08-1 p
8. --- Ühel neist murudest seisab pronksist \VUgras{kuju} \textsuperscript{(28)},
monument ühe meie kaaslinlase mälestuseks, kes tegi oma sünnilinnale suuri teeneid.

%%K t11-tit-3 sub
\VUsaut{4.0mm}
\VUpk{10.2pt}{40}{Sõiduvahendid.}
\VUsaut{3.4mm}

%%K t11-c1-09-1 p
9. --- Seni ei ole meie linn veel elektrivalgusega valgustatud, tal on ainult
\VUgras{gaasipõletid} \textsuperscript{(29)}. Kuid \VUgras{kohalik ajakirjandus} peab
kampaaniat, et seda valgustussüsteemi parandataks, \VUgras{nagu} on juba
täiustatud \VUgras{trammide veosüsteemi}. \VUgras{Vanad} \VUgras{hobuste} \VUgras{veetavad}
vankrid on praktiliselt asendatud \VUgras{elektritrammidega} \textsuperscript{(31)}, mis
veavad reisijaid kiiresti \VUgras{linna} ühest otsast \VUgras{teise} väga \VUgras{odava}
hinna eest.

%%K t11-c1-10-1 p
10. --- Kohtumaja lähedal on \VUgras{pank} \textsuperscript{(32)}, kus diskonteeritakse
kaubandusväärtpabereid. \VUgras{Panga inkassaatorid} \textsuperscript{(33)} käivad kodus
võlgu sisse nõudmas.

%%K t11-tit-4 sub
\VUsaut{4.0mm}
\VUpk{10.2pt}{40}{Kunst ja haridus.}
\VUsaut{3.4mm}

%%K t11-c1-11-1 p
11. --- Ilus hoone, mis seisab lähedal, on \VUgras{muuseum}. See sisaldab koos linna
\VUgras{raamatukogu} \textsuperscript{(34)}, ilusaid \VUgras{loodusteaduslikke kogusid}
\textsuperscript{(35)} (loodusmuuseum) ja graatsilist \VUgras{kunstimuuseumi}
\textsuperscript{(36)}, mis moodustab suurepärase alalise \VUgras{näituse}.

%%K t11-c1-11-2 p
See avatakse üldsusele kaks korda nädalas, kuid võõrad võivad seda külastada igal
päeval väikese jootraha eest \VUgras{valvurile} \textsuperscript{(37)}.
\VUgras{Kunstnikud} \textsuperscript{(38)} tulevad sinna töötama. Neid nähakse
\VUgras{ees} oma molberti, \VUgras{palett} käes, kuulsaid maale kopeerimas.

%%K t11-c1-12-1 p
12. --- Kõrgendikul, muuseumi taga, hakkab paistma \VUgras{kuppel}
\textsuperscript{(40)}, mis kroonib \VUgras{haiglat} \textsuperscript{(39)}, ja
\VUgras{seminari} \textsuperscript{(41)} hooned, kus koolitati meie vallakoolide
õpetajaid.

Teatri platsil on \VUgras{postkontor} \textsuperscript{(43)}, mis on sisse seatud
\VUgras{eramajja} \textsuperscript{(42)}.

%%K t11-tit-5 sub
\VUsaut{5.0mm}
\VUpk{11.4pt}{40}{Tulekahju.}
\VUsaut{3.0mm}

%%K t11-tit-6 sub
\VUpk{10.2pt}{40}{Tulekahju hüüd.}
\VUsaut{3.4mm}

%%K t11-c2-01-1 p
1. --- Hirmuäratav kuuldus levib läbi linna: teatris on äsja puhkenud tulekahju!
Hetkel, mil ma jõuan \VUgras{platsile} \textsuperscript{(96)}, on rahvahulk juba
kogunenud \VUgras{kõnniteele} \textsuperscript{(46)} ja \VUgras{sõiduteele}
\textsuperscript{(47)}. \VUgras{Postiametnikud} \textsuperscript{(44)} ja
\VUgras{kirjakandjad} \textsuperscript{(45)} väljuvad postkontorist. \VUgras{Trammijuht}
\textsuperscript{(48)} pidi oma vaguni peatama, et lasta mööda linna
\VUgras{kiirabivankrit} \textsuperscript{(50)}, mis saabub \VUgras{kirurgi}
\textsuperscript{(52)} ja \VUgras{sanitariga} \textsuperscript{(51)}, et hooldada
\VUgras{haavatut} \textsuperscript{(53)}, kes on toodud \VUgras{kanderaamil}
\textsuperscript{(54)} \VUgras{ajalehekioski} \textsuperscript{(55)} juurde.

%%K t11-c2-02-1 p
2. --- Jalaväekompanii, mis oli õppuselt tagasi tulemas, peatus, et koos linna
\VUgras{ratsapolitseinikega} \textsuperscript{(61)} korrateenistust tagada.
\VUgras{Ratsanikud}, kelle hobused tagajalgadele tõusevad, oskavad paremini kui
\VUgras{sõdurid} \textsuperscript{(57)} või \VUgras{sandarmid} \textsuperscript{(63)}
tagasi ajada \VUgras{uudishimulikke} \textsuperscript{(56)}. Pealegi on sandarmid väga
hõivatud. Üks neist näitab \VUgras{politseinikele} \textsuperscript{(64)}, kuhu tuleb
transportida teine \VUgras{õnnetu} \textsuperscript{(65)}, julge tuletõrjuja, kes kukkus
redelilt.

%%K t11-c2-03-1 p
3. --- \VUgras{Ohvitser} \textsuperscript{(66)} täpsustab koha, kuhu tuleb panna saabuv
\VUgras{aurupump} \textsuperscript{(67)}. Seda võimsat masinat oodates on ta lasknud
kiiresti üles seada \VUgras{käsipumba} \textsuperscript{(68)}, mille pikk \VUgras{voolik}
\textsuperscript{(69)} paiskab vett joana tulele.

Kuus tuletõrjujat käsitsevad seda, sellal kui nende kaaslane, kes hoiab \VUgras{pihustit}
\textsuperscript{(70)}, suunab \VUgras{veejoa} \textsuperscript{(71)} tulekahju keskmesse.

%%K t11-c2-04-1 p
4. --- Teatri teisel küljel on toetatud suur \VUgras{redel} \textsuperscript{(72)}. See
võimaldab tuletõrjujatel läheneda leekidele, mis pritsivad musta \VUgras{suitsu}
\textsuperscript{(75)} keeriste vahelt.

%%K t11-tit-7 sub
\VUsaut{5.0mm}
\VUpk{10.2pt}{40}{Vahvad teod.}
\VUsaut{3.4mm}

%%K t11-c2-05-1 p
5. --- \VUgras{Tulekahju} \textsuperscript{(74)} sai alguse \VUgras{teatri}
\textsuperscript{(73)} fuajees, sellal kui harjutati \VUgras{ooperit}, mille esimene
\VUgras{etendus} pidi toimuma homme. Õnneks oli saalis vähe \VUgras{pealtvaatajaid}
\textsuperscript{(76)}. Vaesed varjusid \VUgras{rõdule} \textsuperscript{(77)}, kuhu
julged \VUgras{tuletõrjujad} \textsuperscript{(86)} tulevad neile appi redeli ja
köitega.

%%K t11-c2-06-1 p
6. --- \VUgras{Sammastiku} \textsuperscript{(78)} all, kus \VUgras{kuulutus}
\textsuperscript{(79)} teatab etendusest, nähakse põgenemas \VUgras{muusikuid}
\textsuperscript{(81)}, kes kuuluvad \VUgras{orkestrisse} \textsuperscript{(80)}, kaasas
oma pillid: \VUgras{viiul} \textsuperscript{(81)}, \VUgras{flööt} \textsuperscript{(82)},
\VUgras{tšello} \textsuperscript{(83)}, \VUgras{tromboon} \textsuperscript{(84)},
\VUgras{trumm} \textsuperscript{(85)} jne. Püütakse päästa osa \VUgras{mööblist}
\textsuperscript{(87)}.

%%K t11-c2-07-1 p
7. --- \VUgras{Näitlejad} \textsuperscript{(89)} ja \VUgras{näitlejannad}
\textsuperscript{(88)}, \VUgras{lauljad} ja \VUgras{lauljannad} pidid põgenema oma
teatri\VUgras{kostüümides} \textsuperscript{(90)}. Tenor seletab
\VUgras{ajakirjanikule} \textsuperscript{(92)}, kuidas see juhtus. \VUgras{Reporter}
tuli kohe esimesest hetkest koos võimuesindajatega: \VUgras{prefekt}
\textsuperscript{(91)}, \VUgras{linnapea} \textsuperscript{(93)},
\VUgras{politseikomissar} \textsuperscript{(94)} ja \VUgras{kohtuametnik}
\textsuperscript{(95)}, kes hakkavad uurima, et otsustada vastutuse üle.

\VUsaut{6.0mm}
\VUfilet{22mm}
